% Revised primary mathematical section, independently audited.
% Intermediate LaTeX fragment; theorem-like environments assumed.
% All dependencies on the incoming witness_discharge package are explicit.

\section{Opening a witness at the value determined by its clause}
\label{ao:main}

The preceding witness-discharge calculus asks for three proofs at an
elimination: the paired formula, an equality identifying its fresh witness,
and a formula to be exported after substitution.  It also asks for the
substitution target and the formula before substitution.  The presentation
below removes the equality proof, target, and pre-substitution formula from
that interface.  A fixed probe of the paired formula determines the value
used in the opened clause.  The branch then proves the original outer goal
without changing it.

The fresh numerical parameter is retained in \emph{every} branch.  In
particular, the selected value is allowed to contain that parameter.  This
choice avoids an additional scope restriction on the selector and is
essential to the exact rule stated here.

\subsection{The complete presentation}

Terms are full normal forms of
\[
 t::=x\mid0\mid S(t)\mid I(n,a;u.h),
\]
under the context-closed iterator reductions
\[
 I(0,a;u.h)\longrightarrow a,
 \qquad
 I(Sn,a;u.h)\longrightarrow h[I(n,a;u.h)/u].
\]
Only $u$ is bound in $h$.  Write $N$ for normalization.  All numerical
substitutions avoid capture and are followed by $N$ on numerical subterms.
Normalization leaves the formula constructors unchanged.  The formulas
are equations and $D_{x,y}B$; their ordinary readings are actual equality
and
\[
 (D_{x,y}B)^\circ
     =\exists x\bigl(B^\circ(x,0)\to\forall y\,B^\circ(x,y)\bigr).
\]
Let $H=(0=0)$, $1=S0$, and $2=S(S0)$.

A context is a finite list of clauses $[\vec A;y.B(y)]$, each having
zero or one antecedent.  Its bound variable $y$ is absent from the
antecedent.  Its reading is $\forall y\,B(y)$ for an empty antecedent
list, and $A\to\forall y\,B(y)$ otherwise.  An ordinary assumption $F$
is the unconditional constant clause $[\,;v.F]$.  The initial context
contains exactly the following fixed axiom clause:
\[
                 \mathcal R=[\,;y.y=y].
 \tag{Ax}\label{ao:axiom}
\]
This explicitly accounts for reflexivity.  The designated contradiction is
$0=S0$.

For a matrix $B(x,y)$ and a fresh parameter $b$, compute
\[
 B_0=N(B(b,0)),\qquad B_1=N(B(b,1)),
\]
and define
\[
 \rho(B,b)=
 \begin{cases}
 t,& B_0\equiv(0=0)\text{ and }B_1\equiv(b=t),\\
 b,&\text{otherwise}.
 \end{cases}
 \tag{Selector}\label{ao:selector}
\]
Here $\equiv$ denotes literal equality of canonical syntax.  The left
side in the successful probe must be exactly the fresh parameter $b$.
The right side $t$ may contain $b$.  No theorem search, equality decision
up to provability, certificate choice, or additional equality proof is
hidden in the selector.

For any well-scoped term $r$ of the branch, abbreviate the canonical clause
\[
 \mathcal C_B(r)=[[N(B(r,0))];y.N(B(r,y))].
\]
Substituting $r$ into this clause is capture avoiding and simultaneous.
An occurrence of $b$ inside $r$ is not substituted again.

The proposed presentation $\mathsf{AO}$ has exactly three instruction
families.
\begin{enumerate}
\item \textbf{Use.} For a retained clause $[\vec A;y.B]\in\Gamma$,
prove its one antecedent when present and infer $N(B(t))$, at any
well-scoped numerical term $t$.  An unconditional clause has no proof
premise.
\item \textbf{Introduce.} From
\[
       \Gamma,B(t,\eta)\vdash B(t,Sz)
       \quad\text{infer}\quad
       \Gamma\vdash D_{x,y}B,
 \tag{I}\label{ao:intro}
\]
where $z$ is fresh.  Set $\eta=z$ exactly when $B$ has equation root and
the canonical witness $t$ is literally $0$; otherwise set $\eta=0$.
The displayed branch assumption is an unconditional constant clause.
\item \textbf{Open.} With $b$ fresh, infer
\[
 \frac{\Gamma\vdash D_{x,y}B\qquad
       \Gamma,\mathcal C_B(\rho(B,b))\vdash A}
      {\Gamma\vdash A}.
 \tag{O}\label{ao:open}
\]
The branch retains the fresh parameter $b$, whether or not the selector
returns it.  The conclusion $A$, the retained context $\Gamma$, and the
passive parameters of the major formula are independent of $b$.
\end{enumerate}
The branch of O proves the same outer goal $A$.  There is no separate
formula $F$ to choose, no substitution in the exported conclusion, no
equality-proof premise, and no supplied replacement term.

These are all the primitive inferences.  Full normalization is part of the
trusted syntax and checker.  There is no additional equality-transport,
computation, numerical-case, induction, weakening, cut, or ordinary paired-
elimination rule.  The axiomatic source of reflexivity remains the fixed
clause~\eqref{ao:axiom}.

\subsection{Direct soundness, including self-containing targets}

\begin{lemma}[The original opened clause determines the selected equality]
\label{ao:selected-equality}
In the context $\Gamma,\mathcal C_B(b)$, the equation
$b=\rho(B,b)$ is derivable by Use and the fixed reflexivity clause.
\end{lemma}
\begin{proof}
If the selector succeeds, the clause's antecedent is literally $H$.
Use $\mathcal R$ at $0$ and then use $\mathcal C_B(b)$ at $1$; this
gives $b=t$.  If the selector fails, use $\mathcal R$ at $b$.
\end{proof}

\begin{theorem}[Soundness of automatic opening]
\label{ao:sound}
Rule O is sound in intuitionistic first-order equality logic over the
ordinary scalar arithmetic, even when $\rho(B,b)$ contains $b$.
\end{theorem}
\begin{proof}
Eliminate the existential in the ordinary reading of the major formula
at a fresh parameter $b$.  Its actual local assumption is
$\mathcal C_B(b)^\circ$.  Lemma~\ref{ao:selected-equality} supplies
$b=\rho(B,b)$.  Ordinary equality substitution in the fixed clause
template $\mathcal C_B(v)^\circ$ gives
$\mathcal C_B(\rho(B,b))^\circ$.  The supplied branch then proves $A$.
Since $A$ and the retained outer assumptions are independent of $b$,
ordinary existential elimination discharges the original witness.

No occurs check is needed on the selected target.  Equality substitution
from $b=t(b)$ to a fixed formula instance $C(t(b))$ is valid; it is a
single simultaneous substitution.  The fresh $b$ remains in the branch,
including in the selected clause when necessary, and the final goal does
not contain it.  The rule does not assert that the equation $b=t(b)$ has
a solution independently of the major premise.
\end{proof}

\subsection{The bootstrap must precede ordinary elimination}

The new branch assumption is generally
$\mathcal C_B(\rho(B,b))$, whereas ordinary paired elimination expects
$\mathcal C_B(b)$.  Thus ordinary elimination cannot be justified merely
by observing that its goal is $b$-free.  We first establish the numerical
substitution theorem needed to change its branch assumption.  This order
avoids assuming the desired equality or elimination principle in its own
proof.

Weakening, injective fresh renaming, and cut of an ordinary assumption are
admissible by direct induction on derivations.  For O, these operations
leave its matrix and determined target unchanged up to fresh renaming.
For cut, replace uses of the unconditional constant clause for the cut
assumption by its supplied proof; freshen local parameters before insertion.
No arbitrary numerical substitution theorem is used in this argument.

\begin{lemma}[Openings on which the selector necessarily falls back]
\label{ao:safe}
The selector returns its fresh parameter in each of the following cases:
\begin{enumerate}
\item the matrix $B$ does not use its witness coordinate $x$;
\item the matrix $B$ has outer constructor $D$.
\end{enumerate}
Consequently, O is literally ordinary paired elimination on these matrices.
\end{lemma}
\begin{proof}
In the first case the normalized probe at $1$ does not contain the fresh
parameter $b$, so it cannot be an equation whose left side is $b$.
In the second case the probe is still a $D$ formula, because numerical
normalization never removes a formula constructor.  Neither can satisfy
the successful guard in~\eqref{ao:selector}.
\end{proof}

For unused $w$ write $q_yA(y)=D_{w,y}A(y)$.  Its eliminations fall under
the first case of Lemma~\ref{ao:safe}.  Therefore, before any general
equality reasoning, we have
\[
 \frac{\Gamma,A(0)\vdash A(Sz)}{\Gamma\vdash q_yA(y)},
 \qquad
 \frac{\Gamma\vdash q_yA(y)\qquad\Gamma\vdash A(0)}
      {\Gamma\vdash A(t)}.
 \tag{Q}\label{ao:q}
\]
The first rule uses witness $1$ in I, which forces case mode.  The second
uses O and then Use on the opened clause $[A(0);y.A(y)]$.  As a consequence,
proofs of $A(0)$ and $A(Sz)$ for fresh $z$ yield a proof of $A(t)$ at every
well-scoped term $t$.  These numerical cases use only fallback openings.

\begin{lemma}[All case introductions]
\label{ao:all-case}
For every matrix $B$ and every canonical witness $t$, the inference
\[
       \Gamma,B(t,0)\vdash B(t,Sz)
       \quad\Longrightarrow\quad
       \Gamma\vdash D_{x,y}B
\]
is derivable in $\mathsf{AO}$.  Its derivation uses only fallback
openings, in addition to Use and I.
\end{lemma}
\begin{proof}
Only the case of an equation matrix $E(x,y)$ and witness $t\equiv0$
is missing from the primitive rule.  Define the normalizing selector
$C(v,a,c)=N(I(v,a;u.c))$, with $u$ absent from $c$, and the equation gate
\[
 G_v(s=t)=(C(v,0,s)=C(v,0,t)),\qquad L(E)=q_vG_vE.
\]
The literal normalized endpoints are $G_0E\equiv H$ and
$G_{Sw}E\equiv E$.  The rules~\eqref{ao:q} and reflexivity therefore
give uniform derivations $E\vdash L(E)$ and $L(E)\vdash E$.  Every
opening in these derivations has an unused witness coordinate.

Suppose the given derivation is
$\Gamma,E(0,0)\vdash E(0,Sz)$.  Set $A(x,y)=L(E(x,y))$.  Under
$A(0,0)$, unwrap the equation, insert the given derivation by ordinary-
assumption cut, and wrap the resulting equation.  This gives $A(0,Sz)$.
Since $A$ has a $D$ root, I at witness $0$ uses cases and gives
$D_{x,y}A(x,y)$.

Open this formula at a fresh $c$.  Its matrix is compound, so this is
a fallback opening by Lemma~\ref{ao:safe}.  It supplies the ordinary
clause $[A(c,0);y.A(c,y)]$.  To introduce $D_{x,y}E(x,y)$ at witness
$c$, assume $E(c,0)$, wrap it as $A(c,0)$, use that clause at $Sz$,
and unwrap $A(c,Sz)$ as $E(c,Sz)$.  The witness $c$ is not literal
zero, so I uses cases.  Close I and the fallback opening.  The result
is the required $D_{x,y}E$.  Neither general equality transport nor
general numerical substitution in derivations has been used.
\end{proof}

\subsection{A full numerical substitution theorem}

\begin{theorem}[Normalized substitution in derivations]
\label{ao:substitution}
For every finite derivation $\Gamma\vdash_{\mathsf{AO}}A$ and every
well-scoped simultaneous substitution $\sigma$ for its free numerical
parameters, there is an effectively constructed derivation
\[
            N(\Gamma\sigma)\vdash_{\mathsf{AO}}N(A\sigma).
\]
The fixed clause $\mathcal R$ is unchanged.  The theorem permits a source
parameter to be mapped to a term containing a receiving parameter with
the same printed name; substitution is simultaneous, not iterative.
\end{theorem}
\begin{proof}
Induct on the finite derivation, with the substitution universally
quantified in the induction hypothesis.  Freshen all bound and local
parameters before applying a substitution.  Use is immediate from
normalization--substitution coherence.  Reflexivity uses the same fixed
axiom clause.

For I, apply the induction hypothesis to its strictly smaller branch,
extending $\sigma$ by a fresh parameter for its local variable $z$.
If the mode is preserved, the resulting rule instance is direct.  An
old progressive instance has a literal zero witness, which stays zero;
the equation root also stays an equation.  The only possible change is
from cases to progress when a previously nonzero witness normalizes to
zero and the matrix is an equation.  In this event use the already proved
all-case rule, Lemma~\ref{ao:all-case}, on the substituted branch.
This is an insertion of a known derived inference, not an appeal to the
substitution theorem on a larger newly constructed proof.

It remains to handle O.  To distinguish source and receiving scopes,
call its old fresh parameter $b$ and its new fresh parameter $c$.
Translate its strictly smaller major proof under $\sigma$, and compute
\[
  B'=N(B\sigma),\qquad r'=\rho(B',c).
\]
There are two cases according to the \emph{old} guard.

If the old guard succeeded, write its old target as $t(b)$.  Apply the
induction hypothesis to the old branch with $\sigma$ on outer parameters
and $b\mapsto c$.  Success is preserved under this protected substitution:
the old base $H$ stays $H$, and the old probe $b=t(b)$ becomes
$c=N(t\sigma[b\mapsto c])$.  Thus the new guard also succeeds and
\[
 r'=N(t\sigma[b\mapsto c]).
\]
By coherence, the substituted branch clause is exactly
$\mathcal C_{B'}(r')$.  This argument does not require $t$ to be
independent of $b$.

If the old guard failed, its old branch clause is $\mathcal C_B(b)$.
Apply the induction hypothesis to that old branch with the substitution
that extends the outer $\sigma$ by
\[
                         b\longmapsto r'.
\]
This is allowed whether $r'=c$ or $r'$ contains $c$ in a larger term.
The source parameter $b$ is replaced in one simultaneous pass.  Its
clause becomes exactly $\mathcal C_{B'}(r')$, again by coherence.
The branch goal and retained outer context did not contain $b$, so
their images are precisely the required $N(A\sigma)$ and
$N(\Gamma\sigma)$.  This covers a guard that stays false and a guard
newly made true by normalization after substitution.

In either case O applies to the translated major and branch, proving
$N(A\sigma)$.  Every recursive call used the induction hypothesis for
a strictly smaller original proof, including when the chosen replacement
contains the new fresh parameter.  The construction is therefore well
founded and yields a finite proof.
\end{proof}

The selector's successful condition is preserved under protected passive-
parameter substitution, but its failed condition need not be.  The second
case of this proof is necessary: treating the selector as an invariant
Boolean test would leave a gap.

\begin{corollary}[Ordinary paired elimination]
\label{ao:ordinary-open}
For arbitrary matrices, not just those of Lemma~\ref{ao:safe}, the rule
\[
 \frac{\Gamma\vdash D_{x,y}B\qquad
       \Gamma,\mathcal C_B(b)\vdash A}
      {\Gamma\vdash A},
 \qquad b\notin\operatorname{FV}(\Gamma,A),
\]
is derivable in $\mathsf{AO}$.
\end{corollary}
\begin{proof}
Compute $r=\rho(B,b)$ in the receiving fresh-parameter scope.  Apply
Theorem~\ref{ao:substitution} to the supplied branch, leaving its outer
parameters fixed and replacing its source fresh parameter by $r$.
Its assumption becomes $\mathcal C_B(r)$.  Its goal remains $A$,
because $A$ has no free source parameter.  The receiving scope retains
$b$, so this is legitimate even if $r$ contains $b$.  The new rule O
now applies.  If the selector falls back, this is just a fresh renaming
or the identity instance of the same argument.
\end{proof}

\subsection{Equality transport from a determined witness}

The numerical selector and parity terms
\[
 C(v,a,c)=N(I(v,a;u.c)),\qquad
 P(y)=N(I(y,0;u.C(u,1,0)))
\]
satisfy the literal canonical identities
\[
 C(0,a,c)\equiv a,\quad C(Sv,a,c)\equiv c,
 \quad P(1)\equiv1,\quad P(2)\equiv0.
\]
No range theorem or induction for parity is needed below.

\begin{theorem}[Uncrossed transport]
\label{ao:transport}
For every fixed numerical term template $h(w)$, the inference
\[
 \frac{\Gamma\vdash s=t\qquad\Gamma\vdash N(h(s))=r}
      {\Gamma\vdash N(h(t))=r}
\]
is derivable in $\mathsf{AO}$.
\end{theorem}
\begin{proof}
Choose fresh $x,y$ and use the two-row matrix
\[
 B(x,y)=\left(
 C(y,0,C(P(y),N(h(x)),x))
 =C(y,0,C(P(y),r,t))\right).
 \tag{Rows}\label{ao:rows}
\]
Its relevant normalized instances are
\[
 B(x,0)\equiv H,\qquad
 B(x,1)\equiv(x=t),\qquad
 B(x,2)\equiv(N(h(x))=r).
\]
Introduce $D_{x,y}B$ at witness $s$.  Its successor goal is
\[
 C(P(Sz),N(h(s)),s)=C(P(Sz),r,t).
\]
Use the numerical-case construction~\eqref{ao:q} on
\[
 A(v)=\bigl(C(v,N(h(s)),s)=C(v,r,t)\bigr)
\]
at target $P(Sz)$.  The zero instance is the supplied $N(h(s))=r$,
and every successor instance is the supplied $s=t$.  Choose the case
eigenparameter fresh independently of $z$.  This construction ignores
I's local assumption, so it works in either mode of I, including when
$s$ is literal zero.

Open the resulting $D$ formula.  Its base is $H$ and its fixed probe at
$1$ is $b=t$, so its selected value is the original outer term $t$.
The branch therefore receives $\mathcal C_B(t)$ directly.  Use this
clause at $2$, with the reflexive antecedent $H$.  Its conclusion is
$N(h(t))=r$, which is the unchanged outer goal.  Rule O finishes.
There is no equality-proof branch and no substitution in the exported
formula.  All equality use has been obtained through the determined
witness in the clause.
\end{proof}

The order of these two rows puts the equality identifying the witness at
the smaller fixed probe $1$.  The other row is read by ordinary Use at
$2$.  The naming method is inherited; the new inference determines and
uses its witness automatically.

\begin{corollary}[Full equality]
\label{ao:full-equality}
Symmetry, transitivity, term congruence, crossed term transport, and
transport in every Eq/$D$ formula template are derivable.
\end{corollary}
\begin{proof}
In Theorem~\ref{ao:transport}, take the identity template and the
reflexive minor $s=s$ to derive $t=s$ from $s=t$.  This gives symmetry.
Taking a reflexive minor for a general template gives congruence, and
the identity template together with symmetry gives transitivity.
Symmetry after uncrossed transport gives crossed transport.  Applying
term transport successively to both sides transports an equation template.

For a compound template $D_{x,y}E(v,x,y)$, induct externally on formula
depth.  From $s=t$ and $D_{x,y}E(s,x,y)$, use the derived ordinary
paired elimination at a fresh $c$.  Introduce $D_{x,y}E(t,x,y)$ at
witness $c$, which forces case mode.  Under $E(t,c,0)$, transport
backward at the strictly smaller matrix using $t=s$, obtaining
$E(s,c,0)$.  Use the old opened clause at $Sz$, then transport the
result forward at the smaller matrix to $E(t,c,Sz)$.  Close I and
ordinary elimination.  This is induction on the finite formula syntax,
not induction inside the arithmetic theory on quantified formulas.
\end{proof}

\subsection{Exact comparison with the predecessor}

Let $\mathsf{WD}$ be the predecessor with the same canonical syntax,
clauses, fixed axiom $\mathcal R$, Use, and I.  Its equality-aware
discharge has a major $D_{x,y}B$, opens $\mathcal C_B(b)$, asks for
branch proofs $b=t$ and $F$, and exports $N(F[t/b])$ when it is
independent of $b$.

\begin{theorem}[The entire consequence relation is unchanged]
\label{ao:equivalence}
For every canonical formula $A$ and every finite canonical context
$\Delta$ of clauses with zero or one antecedent,
\[
 \mathcal R,\Delta\vdash_{\mathsf{AO}}A
 \quad\Longleftrightarrow\quad
 \mathcal R,\Delta\vdash_{\mathsf{WD}}A.
\]
Both directions have effective finite proof transformations, with the
end context and end formula left unchanged.
\end{theorem}
\begin{proof}
First translate a WD proof recursively into AO.  Use and I are shared.
At a WD discharge, the translated original branch proofs give $b=t$
and $F$ under $\mathcal C_B(b)$.  Full formula transport,
Corollary~\ref{ao:full-equality}, gives $N(F[t/b])$ there.  That result
was required by WD to be independent of $b$, so derived ordinary paired
elimination, Corollary~\ref{ao:ordinary-open}, exports it.  The lemmas
used here were proved in AO before this translation and without a
translation of the node currently being replaced.  There is no circular
appeal to an omitted WD rule.

Conversely, translate an AO proof recursively into WD.  Use and I are
again shared.  For its O node, WD has ordinary paired elimination and
all ordinary equality reasoning.  Open its major at $b$, obtaining
$\mathcal C_B(b)$.  If the selector falls back, the translated AO
branch already has the required assumption, so ordinary elimination
finishes directly.

If the selector succeeds with target $t(b)$, Use and reflexivity give
$b=t(b)$.  The base of $\mathcal C_B(b)$ is $H$, so for every fresh
$z$ its clause gives $B(b,Sz)$.  Transport in the witness coordinate
using $b=t(b)$ gives $B(t(b),Sz)$.  Case introduction at the unused
witness $1$ proves the dummy-witness closure
$q_y B(t(b),y)$, ignoring its temporary zero-base assumption.  Ordinary
elimination of this closure supplies exactly the desired clause
$\mathcal C_B(t(b))$.  The recursively translated AO branch can now
be used, weakened by the other retained assumptions and freshened at
the unused dummy witness.  It proves $A$.  Close the dummy-witness
elimination and the original paired elimination.  The final goal is
independent of $b$, as required by O.

This construction remains valid if $t$ contains $b$: throughout the
inner construction $b$ is an available outer parameter.  Only the final
ordinary elimination discharges it.  Each replacement is a finite
derivation built from already established WD rules.  Freshen the finitely
many local parameters before insertion.
\end{proof}

Composing this identity-language equivalence with the predecessor's
canonical/raw comparison and the native/scalar/iPRA bridge yields
bi-interpretability with full first-order intuitionistic primitive recursive
arithmetic.  Both numerical domains are whole, equality is actual equality,
and both numerical composite isomorphisms retain the graph $x=y$.
The two formula round trips are preserved because this comparison changes
neither formula syntax nor ordinary readings.  The preserved contradiction
$0=S0$ and the two effective proof transformations also give equiconsistency.
This conclusion is not inferred from standard-natural-number agreement.

\subsection{What has been simplified and what remains trusted}

Both presentations have three instruction families and the same fixed
reflexivity axiom.  The new O instruction has two proof premises, compared
with the predecessor's three.  It removes the supplied equality target,
the separate equality-proof subtree, and the supplied branch formula
before export.  Its branch proves the unchanged goal.  In exchange, the
kernel computes two fixed matrix instances, tests the literal guard, and
substitutes the determined value into its added clause.  Whole-checker
measurements must count these operations and the unchanged full normalizer.
A smaller certificate interface is a mathematical reduction, but it does
not by itself establish a repository-wide minimum in program size.

The normalizer is a metasyntactic algorithm, not an object-language
universal evaluator.  No uniform primitive-recursive normalization bound
is asserted.  A fuel-bounded certificate checker can charge every
normalization and recursive verification step, including both selector
probes and the selected-clause substitution, to one supplied fuel account.
Each accepted finite proof has some sufficiently large finite fuel value,
so the bounded interface presents the same theorems.  This adds no
arithmetic totality axiom and changes neither interpretation domain.

The new contribution is the determined opening rule and its complete
bootstrap: fallback openings suffice for numerical cases and the
zero-witness repair; these support full substitution of derivations;
substitution recovers ordinary paired elimination; the two-row package
recovers all equality; and the entire consequence relation follows.
The scalar iterator, normal form syntax, paired binder, clause interface,
fixed reflexivity axiom, naming package, and arithmetic bridge are inherited.
The universal arguments above are written proofs, rather than a new
Lean or Rocq verification.  Finite tests of a checker would provide
implementation evidence, not replace either interpretation proof.
